\part[Гармонический анализ.]{Гармонический анализ.}%\protect\linebreak \vbox{\vspace*{5\baselineskip}}\centerline{\includegraphics[height=0.3\textheight]{pictures/Foxi/5}}}

\chapter{Достаточные условия сходимости тригонометрического ряда Фурье в точке.}
\section{Ортогональные системы и ряды Фурье}
\subsection{Ортогональные и ортонормированные системы функций}
Пусть задана система функций 
\begin{equation} \label{ch17.1eq00}
\phi_1(x), \,\phi_2(x), \, \ldots, \, \phi_n(x), \, \ldots,
\end{equation}
области определения которых имеют непустое пересечение. Тогда любой функциональный ряд вида
\begin{equation} \label{ch17.1eq01}
 \sum_{n = 1}^{\infty} a_n \phi_n(x),
 \end{equation}
 где $a_n$ --- некоторые числа, называется \textit{рядом по системе функций \eqref{ch17.1eq00}}, а числа $a_n$ --- \textit{коэффициентами ряда \eqref{ch17.1eq00}.}
 
 Говорят, что \textit{функция $f(x)$ разложена в ряд по системе функций \eqref{ch17.1eq00}}, если указаны числа $a_1, \,a_2, \,\ldots, \, a_n, \,\ldots$, такие, что ряд \eqref{ch17.1eq01} в любой точке области определения функции $f$ сходится и 
 $$
  \sum_{n = 1}^{\infty} a_n \phi_n(x)=f(x) \quad \forall x \in D_f
  $$ 



\begin{defn}
Любые две функции $\phi(x)$	и $\psi(x)$, определенные на промежутке $\Delta$, называются \textit{ортогональными на $\Delta$}, если их произведение интегрируемо на $\Delta$ (в собственном или несобственном смысле) и
$$
\int \limits_\Delta\phi(x) \psi(x)dx=0
$$
Очевидно, функция, тождественно равнвю нулю на промежутке $\Delta$, ортогональна на $\Delta$ любой функции, определенной на $\Delta$.
\end{defn}

\begin{defn}
Система функций, определенных на промежутке $\Delta$, называется \textit{ортогональной на $\Delta$}, если любые две функции этой системы ортогональны на $\Delta$.
\end{defn}

\begin{defn}
Система функций 
\begin{equation} \label{ch17.1eq0}
\{1,\, \cos kx,\, \sin kx\}, \quad k\in \bbN
\end{equation}
называется тригонометрической системой функций.			
\end{defn}
Легко проверить, что тригонометрическая система \eqref{ch17.1eq0} ортогональна на интервале $(-\pi;\pi)$. А так как все функции это системы периодические с периодом $2\pi$, то справедливо следующее утверждение:

\textit{Тригонометрическая система \eqref{ch17.1eq0} ортогональна на любом промежутке длины $2\pi$}.	

Обобщение системы \eqref{ch17.1eq0} является система функций
$$
1,\ \cos \frac{\pi}{l}x,\ \sin \frac{\pi}{l}x,\ldots,\cos n\frac{\pi}{l}x,\ \sin n\frac{\pi}{l}x,\ldots,
$$
где $l>0$, которую тоже будем называть \textit{тригонометрической системой}. Очевидно, она ортогональная на любом промежутке длины $2l$.
\begin{defn}
Система функций \eqref{ch17.1eq0}, определенных на промежутке $\Delta$, называется \textit{ортонормированной на $\Delta$}, если она ортогональна на $\Delta$ и 
\begin{equation} \label{ch17.1eq000}
\int \limits_\Delta{|\phi_n(x)|^2dx=1} \quad \forall n \in \bbN.
\end{equation}
Условие \eqref{ch17.1eq000} называется \textit{условием нормировки.}
\end{defn}

\subsection{Ряды Фурье по ортогональным системам функций}

Пусть $f(x),\ x \in \Delta$, разложена в ряд по системе функций
\begin{equation} \label{ch17.1eq02}
\phi_1(x), \,\phi_2(x), \, \ldots \phi_n(x), \,\ldots,
\end{equation}
ортогональной на промежутке $\Delta$:
\begin{equation} \label{ch17.1eq03}
f(x)= \sum_{n = 1}^{\infty} a_n \phi_n(x),
\end{equation}
Найдем формулы для определения коэффициентов этого ряда. Для этого обе части равентсва \eqref{ch17.1eq03} умножим на $\phi_k(x)$ и проинтегрируем по промежутку $\Delta$, предполагая, что полученный ряд можно интегрировать почленно. В результате получим равенство
$$
\int \limits_\Delta{f(x) \phi_k(x)}dx=a_k\int \limits_\Delta|\phi_k(x)|^2dx,
$$
так как интегралы от произведений $\phi_n(x)\psi_n(x)$ при $n\ne k$ равны нулю (в силу ортогональьности системы \eqref{ch17.1eq02})

Предположим еще, что ортогональную систему \eqref{ch17.1eq02} можно нормировать, т.е. что
$$
\int \limits_\Delta|\phi_k(x)|^2dx\ne 0 \quad \forall n
$$
Тогда для коэффициентов $a_n$ ряда \eqref{ch17.1eq03} получим формулу
\begin{equation} \label{ch17.1eq04}
a_n= \frac{1}{\int \limits_\Delta|\phi_n(x)|^2dx} \int \limits_\Delta f(x)\phi_n(x)dx, \quad n\in \bbN.
\end{equation}

\begin{defn}
Числа $a_n$, определяемые по формулам \eqref{ch17.1eq04} называются \textit{коэффициентами Фурье функции $f$ по ортогональной системе \eqref{ch17.1eq0}}, а ряд $\sum\limits_{n = 1}^{\infty} a_n \phi_n(x)$ называется \textit{рядом Фурье функции $f$ по этой системе.} 	
\end{defn}

\begin{defn} 
Функция $f$, определенная на конечном промежутке $\Delta=(a;b)\subset \bbR$, за исключением, быть может, конечного числа точек из $\Delta$, называется \textit{абсолютно интегрируемой} на промежутке $\Delta$, если существует конечный набор элементов $c_0, \,c_1,\,\ldots,\,c_m \in \bboR$, такой, что
\begin{enumerate}

\item[1)]	
$a=c_0<c_1<\ldots<c_m=b$;
\item[2)]
$\forall [\alpha;\beta]\subset(c_{k-1};c_k),\ k\in \overline{1,m}$ функция $f$ интегрируема по Риману на $[\alpha;\beta]$;	
\item[3)]
$\int \limits_\Delta |f(x)|dx$ сходится;
\item[2)]
$ \forall\epsilon>0$ функция $f$ интегрируема по Риману на множестве 

$G_{\epsilon}=\Delta \setminus (\bigcup \limits_{k=0}^{m}(O_{\epsilon}(c_k))$
\end{enumerate}	
\end{defn}
\section{Тригонометрические ряды Фурье}	
	
Для любой функции $f$, определенной и абсолютно интегрируемой на конечном интервале $(a;b)$, определены коэффициенты Фурье по тригонометрической системе, ортогональной на интервале $(a;b)$:
\begin{equation} \label{ch17eq1}
1,\ \cos \frac{\pi}{l}x,\ \sin \frac{\pi}{l}x,\ldots,\cos n\frac{\pi}{l}x,\ \sin n\frac{\pi}{l}x,\ldots,
\end{equation}
где $l = (b - a)/2$. Соответствующий ряд Фурье обычно записывают в следующем виде:
\begin{equation} \label{ch17eq2}
\frac{a_0}{2} + \sum_{n = 1}^{\infty} a_{n}\,\cos n\frac{\pi}{l}x + b_{n}\, \sin n\frac{\pi}{l}x
\end{equation}
Из общей формулы для коэффициентов Фурье следует, что коэффициенты $a_0, a_n, b_n$ ряда ($\ref{ch17eq2}$) вычисляются по формулам:
\begin{align} \label{ch17eq3}
a_0 &= \frac{1}{l} \int\limits_{a}^{b} f(x)dx,\\
a_n &= \frac{1}{l} \int\limits_{a}^{b} f(x)\,\cos n\frac{\pi}{l}x\, dx,\\
b_n &= \frac{1}{l} \int\limits_{a}^{b} f(x)\,\sin n\frac{\pi}{l}x\, dx,
\end{align}

\begin{defn}
Числа $a_0, a_n, b_n$, определяемые по формулам выше, называются \textit{коэффициентами Фурье функции f по тригонометрической системе} \eqref{ch17eq1}.
\end{defn}

Так как общий случай заменой $\xi = \cfrac{\pi}{l}\, x$ сводится к тригонометрической системе, ортогональной на $(-\pi;\pi)$, то в дальнейшем будем изучать в основном только ряды по тригонометрической системе
$$
1,\ \cos x,\ \sin x,\ldots,\cos nx,\ \text{sin}\, nx,\ldots,
$$
которую будем называть \textit{стандартной тригонометрической системой.}

Вместо тригонометрической системы \eqref{ch17eq1} ортогональной на интервале $(-l;l)$, часто рассматривают систему комплекснозначных функций
\begin{equation} \label{ch17eq6}
\phi_{\nu}(x) = e^{i\nu\frac{\pi}{l}x},\quad \nu \in \bbZ,
\end{equation}
Ее тоже называют тригонометрической, так как
$$
e^{i\nu\frac{\pi}{l}x} = \cos \nu\frac{\pi}{l}x + i\, \sin \nu\frac{\pi}{l}x
$$
\begin{defn}
Для любой функции $f(x) \in \overset{*}{L^{R}_1}(-l;l)$ числа
\begin{equation}
c_{\nu} = \frac{1}{2l} \int\limits_{-l}^{l} f(x)\, e^{-i\nu\frac{\pi}{l}x}dx,\quad \nu = 0,\pm 1, \pm 2,\ldots
\end{equation}
называются \textit{коэффициентами Фурье функции f по ортогональной системе} \eqref{ch17eq6}, причем
\begin{equation}
c_0 = a_0/2, \quad c_{k} = \cfrac{a_k - ib_k}{2}, \quad c_{-k} = \cfrac{a_k + ib_k}{2} \quad \forall\ k \in \bbN,
\end{equation}
где $a_0, a_n, b_n$ "--- коэффициенты Фурье функции $f$ по тригонометрической системе \eqref{ch17eq1}.
\end{defn}

\begin{defn}
Выражение
$$
\sum_{\nu = -\infty}^{\infty} c_{\nu}e^{i\nu\frac{\pi}{l}x},
$$
где $c_{\nu} = c_{\nu}(f)$, называют \textit{рядом Фурье функции f по системе} \eqref{ch17eq6}, а сумма
$$
T_{n}(f;x) = \sum_{\nu = -n}^{n} c_{\nu}e^{i\nu\frac{\pi}{l}x}
$$
называется \textit{n-й частичной суммой} этого ряда.
\end{defn}

\section{Теорема Римана об осцилляции}
\begin{defn}
Для любой функции $f$ замыкание множества точек $x \in D_{f}$, в которых $f(x) \neq 0$, называется \textit{носителем функции f} и обозначается $\supp f$.
\end{defn}
\begin{defn}
Функция, определенная на $\bbR$, называется \textit{финитной\rindex{функция!финитная}}, если ее носитель ограничен, т.е. если она равна нулю вне некоторого отрезка.
\end{defn}
\begin{defn}
Функция, определенная на некотором промежутке $\Delta$, называется \textit{ступенчатой\rindex{функция!ступенчатая}}, если существует разбиение промежутка $\Delta$ на конечное число промежутков, на каждом из которых она постоянна.
\end{defn}
\begin{thm} \label{ch17thm1}
Если функция $f(x)$ абсолютно интегрируема на промежутке $\Delta$, то для любого $\epsilon > 0$ существует финитная ступенчатая функция $\phi(x)$ такая, что $\supp\phi \subset \overline{\Delta}$ и
\begin{equation} \label{ch17eq7}
\int\limits_{\Delta} |f(x) - \phi(x)|dx < \epsilon.
\end{equation}
\end{thm}
\begin{proof}
Выберем некоторое $\epsilon > 0$ и построим финитную ступенчатую функцию $\phi(x)$, удовлетворяющую неравенству.

Из определения абсолютно интегрируемой функции $f$ на промежутке $\Delta$ следует, что существует ограниченное измеримое множество $g_{\epsilon} \subset \Delta$, на котором функция $f$ интегрируема по Риману, причем
$$
\int\limits_{\Delta} |f(x)|dx - \int\limits_{g_{\epsilon}} |f(x)|dx = \int\limits_{\Delta\backslash g_{\epsilon}} |f(x)|dx
$$

Через $f_{\epsilon}(x)$ обозначим функцию, равную $f(x)$ на $g_{\epsilon}$ и нулю вне $g_{\epsilon}$. Очевидно,
\begin{equation} \label{ch17eq8}
\int\limits_{\Delta} |f(x) - f_{\epsilon}(x)|dx = \int\limits_{\Delta\backslash g_{\epsilon}} |f(x)|dx < \frac{\epsilon}{2}
\end{equation}

Функция $f_{\epsilon}(x)$ равна нулю вне ограниченного множества $g_{\epsilon} \subset \Delta$, поэтому ее носитель содержится в некотором отрезке $[a;b] \subset \overline{\Delta}$. На этом отрезке она интегрируема по Риману, так как она интегрируема по Риману на $g_{\epsilon}\ \text{и на}\ [a;b]\backslash g_{\epsilon}$. Поэтому существует разбиение $\tau$ отрезка $[a;b]$ на промежутки $\Delta_1,\ldots,\Delta_N$ такое, что
$$
\int\limits_{a}^{b} f_{\epsilon}(x)dx - s(f_{\epsilon};\tau) < \frac{\epsilon}{2},
$$
где $s(f_{\epsilon};\tau)$ "--- нижняя сумма Дарбу функции $f_{\epsilon}(x)$, т.е.
$$
s(f_{\epsilon};\tau) = \sum_{j = 1}^{N} m_{j}|\Delta_j|,
$$
где $m_j = \inf\limits_{x \in \Delta_j}\,f_{\epsilon}(x)$. Через $\phi(x)$ обозначим ступенчатую функцию, которая равна $m_j$ на $\Delta_j,\ j = 1,\ldots,N$, и нулю вне $[a;b]$. Очевидно, $\phi(x) \le f_{\epsilon}(x)$ для любого $x \in \bbR$,
$$
\int\limits_{\Delta} \phi(x)dx = \int\limits_{a}^{b} \phi(x)dx = s(f_{\epsilon};\tau),
$$
и, следовательно,
\begin{equation} \label{ch17eq9}
\int\limits_{\Delta} |f(x) - \phi(x)|dx = \int\limits_{a}^{b} f_{\epsilon}(x)dx - s(f_{\epsilon};\tau) < \frac{\epsilon}{2}
\end{equation}
Из неравенств \eqref{ch17eq8} и \eqref{ch17eq9} следует, что финитная ступенчатая функция $\phi(x)$ удовлетворяет неравенству \eqref{ch17eq7} теоремы. Действительно:
\begin{multline*}
\int\limits_{\Delta} |f(x) - \phi(x)|dx \le \int\limits_{\Delta} |f(x) - f_{\epsilon}(x)|dx+\\
+\int\limits_{\Delta} |f(x) - \phi(x)|dx < \frac{\epsilon}{2} + \frac{\epsilon}{2} = \epsilon
\end{multline*}
\end{proof}

\begin{thm}[Римана об осцилляции] \label{ch17thm2}
Если функция $f(x)$ абсолютно интегрируема на промежутке $\Delta$, то
\begin{equation} \label{ch17eq10}
\lim_{\lambda \to \infty}\int\limits_{\Delta} f(x)e^{i\lambda x}dx = 0
\end{equation}
\end{thm}
\begin{proof}
Для характеричстических функций любого конечного промежутка это утверждение очевидно. Действительно, если $\xi$ и $\eta$ "--- концы промежутка, то
$$
\left|\int\limits_{\xi}^{\eta} e^{i\lambda x}dx\right| = \left|\cfrac{e^{i\lambda \eta} - e^{i\lambda \xi}}{i\lambda}\right| \le \frac{2}{|\lambda|}.
$$

А так как любая финитная ступенчатая функция есть линейная комбинация характеристических функций конечного числа конечных промежутков, то утверждение теоремы верно для любой такой функции.

Согласно теореме \ref{ch17thm1}, для любого $\epsilon > 0$ существует финитная ступенчатая функция $\phi(x)$ такая, что
$$
\int\limits_{\Delta} |f(x) - \phi(x)|dx < \frac{\epsilon}{2}.
$$

Тогда
$$
\left|\int\limits_{\Delta} f(x)e^{i\lambda x}dx\right| \le \int\limits_{\Delta} |f(x) - \phi(x)|dx\ +\ \left|\int\limits_{\Delta} \phi(x)e^{i\lambda x}dx\right|.
$$

Последний интеграл, согласно уже доказанному, стремится к нулю при $\lambda \to \infty$, поэтому
$$
\exists\ \lambda_{\epsilon} : \quad \forall\ \lambda, |\lambda| > \lambda_{\epsilon} : \left|\int\limits_{\Delta} \phi(x)e^{i\lambda x}dx\right| < \frac{\epsilon}{2}.
$$
Следовательно,
$$
\left|\int\limits_{\Delta} f(x)e^{i\lambda x}dx\right| < \epsilon \quad \forall\ \lambda, |\lambda| > \lambda_{\epsilon},
$$
что и доказывает равенство \eqref{ch17eq10}.
\end{proof}

Отметим, что здесь параметр $\lambda$ может принимать любое значение из $\bbR$ и, в частности, может стремиться как к $+\infty$, так и к $-\infty$. Следовательно, наряду с \eqref{ch17eq10} выполняется и равенство
$$
\lim_{\lambda \to \infty}\int\limits_{a}^{b} f(x)e^{-i\lambda x}dx = 0
$$

\begin{cons}
Если функция $f(x)$ абсолюно интегрируема на промежутке $\Delta$, то
\begin{align*}
&\lim_{\lambda \to \infty}\int\limits_{\Delta} f(x)\cos{\lambda x}\,dx = 0\\
&\lim_{\lambda \to \infty}\int\limits_{\Delta} f(x)\sin{\lambda x}\,dx = 0
\end{align*}
\end{cons}

\section{Признаки сходимости тригонометрических рядов Фурье в точке}
\subsection{Признак Липшица}

Пусть функция $f(x)$ определена в некоторой окрестности точки $x_0$. Говорят, что она в этой точке удовлетворяет $\textit{условию Липшица порядка } \alpha > 0$, если $\exists$ постоянные $C$ и $\delta > 0$ такие, что 

\begin{equation} \label{ch17eq4}
|f(x_0 + \xi) - f(x_0)| \le C |\xi|^\alpha \quad \forall \xi \in (-\delta; \delta).
\end{equation}

\begin{thm} [Признак Липшица сходимости ряда Фурье в точке] \label{ch17thm3}
Если функция $f(x) \in \overset{*}{L^{R}_1}(-\pi;\pi)$ в точке $x_0 \in \bbR$ удовлетворяет условию Липшица порядка $\alpha > 0$, то ее ряд Фурье в точке $x_0$ сходится к $f(x_0)$.
\end{thm}

\begin{proof}
Напомним, что для $n$-й частичной суммы ряда Фурье функции $f$ справедлива формула

$$
T_n(f;x_0) = \frac{1}{2\pi} \int\limits_{-\pi}^{\pi} f(x_0 + \xi) D_n(\xi) \,d\xi,
$$
где

$$
D_n(\xi) = \frac{\sin{\lambda_n \xi}}{\sin{\frac{\xi}{2}}}, \; \frac{1}{2\pi}\int\limits_{-\pi}^{\pi} D_n(\xi) \,d\xi = 1, \; \lambda_n = n + \frac12.
$$

Поэтому 

$$
T_n(f; x_0) - f(x_0) = \frac{1}{2\pi}\int\limits_{-\pi}^{\pi}(f(x_0 + \xi) - f(x_0))D_n(\xi)\,d\xi.
$$

По условию функция $f$ в точке $x_0$ удовлетворяет условию Липшица \eqref{ch17eq4}. Не ограничивая общности, можно считать, что $\delta < \pi$. Тогда функция 

$$
F(\xi) = \frac{f(x_0 + \xi) - f(x_0)}{\sin{\frac{\xi}{2}}}
$$
удовлетворяет неравенству

$$
|F(\xi)| \le \frac{C|\xi|^\alpha}{\left| \sin{\frac{\xi}{2}} \right|} \quad \forall \xi \in (-\delta; \delta).
$$
А так как функция $\sin{\frac{\xi}{2}}$ выпукл вверх на отрезке $[0; \pi]$, то

$$
\sin{\frac{\xi}{2}} \ge \frac{\xi}{\pi} \quad \forall \xi \in [0; \pi],
$$

и поэтому 

$$
|F(\xi)| \le \pi C |\xi|^{\alpha - 1} \quad \forall \xi \in [0; \pi].
$$

Следовательно,
\begin{multline} \label{ch17eq5}
|T_n(f; x_0) - f(x_0)| \le \frac{1}{2} C \int\limits_{-h}^{h} |\xi|^{\alpha - 1}\,d\xi + \\
+ \frac{1}{2\pi} \left| \int\limits_{-\pi}^{-h} F(\xi) \sin{\lambda_n \xi} \,d\xi \right| + \frac{1}{2\pi} \left| \int\limits_{h}^{\pi} F(\xi) \sin{\lambda_n \xi} \,d\xi \right|
\end{multline}
для любого $h \in (0; \delta)$ и любого $n \in \bbN$. Для каждого $h \in (0;\delta)$ функция $F(\xi)$ абсолютно интегрируема на интервалах $(-\pi; -h)$ и $(h; \pi)$. Согласно теореме Римана об осцилляции, интегралы по этим интервалам стремятся к нулю при $n \to \infty$. Поэтому из \eqref{ch17eq5} следует неравенство

$$
\overline{\lim\limits_{n \to \infty}}|T_n(f; x_0) - f(x_0)| \le C \cdot \frac{1}{\alpha}h^{\alpha}.
$$
А так как оно справедливо для любого $h \in (0; \delta)$, то 

$$
\overline{\lim\limits_{n \to \infty}}|T_n(f; x_0) - f(x_0)| = 0.
$$

Очевидно, если верхний предел неотрицательной последовательности равен нулю, то это последовательность сходится и ее предел равен нулю.
\end{proof}

\begin{cons}
Если функция $f(x) \in \overset{*}{L^{R}_1}(-\pi;\pi)$ в точке $x_0 \in \bbR$ дифференцируема или непрерывна и имеет конечные односторонние призводные, то в этой точке ее ряд Фурье сходится к $f(x_0)$.
\end{cons}

\begin{cons}
Если функция $f(x), \: x \in [a;b], \: a < b,$ на интервале $(a;b)$ непрерывна, абсолютно интегрируема и в каждой точке $x \in (a;b)$ дифференцируема или имеет конечные односторонние производные, то в любой точке $x \in (a;b)$ ее ряд Фурье сходится и 

$$
f(x) = \frac{a_0}{2} + \sum\limits_{n = 1}^{\infty} a_n \cos{n \frac{\pi}{l} x} + b_n \sin{n \frac{\pi}{l} x},
$$
где $l = \cfrac{b - a}{2}, \: a_0 = a_0(f), \: a_n = a_n(f), \: b_n = b_n(f)$. Если, кроме того, функция $f$ непрерывна на $[a;b], \: f(a) = f(b)$ и существуют конечные односторонние производные $f'_+(a)$ и $f'_-(b)$, то и в точках $x = a, \: x = b$ ряд Фурье сходится к $f(x)$.
\end{cons}

%TODO Нужны ли односторонние условия Липшица

\subsection{Признак Дини}
Пусть функция $f(x)$ определена в окрестности точки $x_0$ и в этой точке непрерывна. Если, кроме того, существует $\delta > 0$ такое, что разностное отношение
$$
\frac{f(x_0 + \xi) - f(x_0)}{\xi}
$$
абсолютно интегрируемо на интервале $(-\delta;\delta)$, то говорят, что функция $f(x)$ удовлетворяет \textit{условию Дини}.

Очевидно, если функция $f(x)$ в точке $x_0$ удовлетворяет условию Липшица порядка $\alpha > 0$, то в этой точке она удовлетворяет и условию Дини. Обратное утверждение является неверным.

\begin{thm} [Признак Дини сходимости ряда Фурье в точке] \label{ch17thm4}
Если функция $f(x) \in \overset{*}{L^{R}_1}(-\pi;\pi)$ в точке $x_0 \in \bbR$ удовлетворяет условию Дини, то ее ряд Фурье в этой точке сходится к $f(x_0)$.
\end{thm}
\begin{proof}
Как и при доказательстве теоремы \ref{ch17thm3} для $h \in (0;\delta)$ получаем неравенство
$$
\overline{\lim\limits_{n \to \infty}}|T_n(f; x_0) - f(x_0)| \le \frac{1}{2\pi}\int\limits_{-h}^{h}|F(\xi)|d\xi,
$$
где
$$
|F(\xi)| \le \pi\left|\frac{f(x_0 + \xi) - f(x_0)}{\xi}\right|.
$$
Отсюда и из условия Дини следует, что
$$
\overline{\lim\limits_{n \to \infty}}|T_n(f; x_0) - f(x_0)| = 0.
$$
\end{proof}

%TODO Нужны ли односторонние условия Дини

\subsection{Признак Дирихле}
%TODO Тут дофига используется, в том числе односторонние условия и ядро Дирихле. Требует ревизии
Будем говорить, что функция $f(x)$, определенная в некоторой окрестности точки $x_0$, в этой точке удовлетворяет \textit{условию Дирихле}, если существует $\delta > 0$ такое, что на интервалах $(x_0 - \delta;x_0)$ и $(x_0;x_0 + \delta)$ она монотонна и ограничена.

Очевидно, функция $f(x)$, удовлетворяющая условию Дирихле в точке $x_0$, в этой точке имеет конечные односторонние пределы $f(x_0\pm 0)$.
\begin{thm} \label{ch17thm5}
Если функция $f(x) \in \overset{*}{L^{R}_1}(-\pi;\pi)$ в точке $x_0 \in \bbR$ удовлетворяет условию Дирихле, то ее ряд Фурье в этой точке сходится к $M_f(x_0)$.
\end{thm}
\begin{proof}
Как известно,
\begin{multline*}
|T_n(f; x_0) - M_f(x_0)| \le \frac{1}{2\pi} \sum_{\pm} \left|\int\limits_{h}^{\pi} F_{\pm}(\xi) \sin{\lambda_n \xi} \,d\xi \right|\ +\\
+\ \frac{1}{2\pi} \sum_{\pm} \left|\int\limits_{0}^{h} (f(x_0\pm \xi) - f(x_0\pm 0))D_n(\xi) \,d\xi \right|
\end{multline*}
для любого $h \in (0;\pi)$.

Согласно условию Дирихле, существует $\delta > 0$ такое, что функции $f(x_0\pm \xi)$ на интервале $(0;\delta)$ монотонны и ограничены. По второй теореме о среднем для любого $h \in (0;\delta)$ существуют $\theta_{\pm} \in [0;h]$ такие, что
$$
\int\limits_{0}^{h} (f(x_0\pm \xi) - f(x_0\pm 0))D_n(\xi) \,d\xi\ = (f(x_0\pm h) - f(x_0\pm 0))\, \int\limits_{\theta_{\pm}}^{h} D_n(\xi)\,d\xi
$$

Как известно, ядро Дирихле обладает следующим свойством: существует постоянная $C > 0$ такая, что
$$
\left|\int\limits_{\alpha}^{\beta} D_n(\xi)\,d\xi\right| \le 2C\pi
$$
для любых $\alpha, \beta \in (0;\pi)$ и любого $n \in \bbN$. Следовательно для любого $h \in (0;\delta)$ и любого $n \in \bbN$ справедливо неравенство
\begin{multline*}
|T_n(f; x_0) - M_f(x_0)| \le\\
\le C\sum_{\pm}|f(x_0\pm h) - f(x_0\pm 0)| + \frac{1}{2\pi} \sum_{\pm} \left|\int\limits_{h}^{\pi} F_{\pm}(\xi) \sin{\lambda_n \xi} \,d\xi \right|,
\end{multline*}
из которого следует, что
$$
\overline{\lim\limits_{n \to \infty}}|T_n(f; x_0) - M_f(x_0)| \le C\sum_{\pm}|f(x_0\pm h) - f(x_0\pm 0)|
$$
для любого $h \in (0;\delta)$. А так как в последнем неравенстве правая часть стремится к нулю при $h \to +0$, то
$$
\lim\limits_{n \to \infty}|T_n(f; x_0) - M_f(x_0)| = 0.
$$
\end{proof}

\begin{cons}
Если функция $f(x)$ ограничена и кусочно-монотонна на интервале $(-\pi;\pi)$, то в любой точке $x \in (-\pi;\pi)$ ее ряд Фурье сходится к
$$
M_f(x) = \frac{f(x+0) + f(x-0)}{2},
$$
а в точках $-\pi$ и $\pi$ он сходится к $\cfrac{f(-\pi+0) - f(\pi-0)}{2}$. В частности, в точках $x \in (-\pi;\pi)$, где $f$ непрерывна, ряд Фурье сходится к $f(x)$.
\end{cons}
\begin{cons}
Если непрерывная $2\pi$-периодическая функция $f$ кусочно-монотонна на отрезке $[-\pi;\pi]$, то в любой точке $x \in \bbR$ ее ряд Фурье сходится к $f(x)$.
\end{cons}
